\documentclass{scrartcl}
\usepackage[utf8]{inputenc}
\usepackage{fancyhdr}
\usepackage{graphicx}
\usepackage{dirtytalk}
\usepackage[portrait, margin=0.5in, top=0.4in, bottom=0.3in]{geometry}
\usepackage{amsmath, amssymb}
\setlength{\parskip}{1em}
\setlength{\parindent}{0cm} 
\usepackage[mdthm,simplethm]{test}
\usepackage[utf8]{inputenc}
\usepackage{float}
\usepackage{physics}
\setlength{\parindent}{0pt}

\title{Problemset 1}
\author{albert ye}
\date{May 2022}

\makeatletter
\newcommand{\skipitems}[1]{%
  \addtocounter{\@enumctr}{#1}%
}
\makeatother

\begin{document}
\maketitle
\begin{enumerate}
	\item We know that the force of friction must be equivalent to the force of gravity on the mass, because
	the forces of action and reaction must be equal for the mass to remain at rest. Therefore, we see that The
	frictional force must equal $\langle \begin{bmatrix}0 \\ 9.8 \end{bmatrix},
	\begin{bmatrix} \cos 30 \\ \sin 30 \end{bmatrix} \rangle = 9.8 \sin 30 = \boxed{4.9}$.
	\skipitems{1}
	\item Because $\vec{AB} + \vec{BC} = \vec{AC}$, we can extend this to get $\vec{AB} + \vec{BC} + \cdots +
	\vec{DE} = \vec{AE}$. Therefore, the expression evaluates to $\vec{AE} + \vec{EA} = \vec{AE} - \vec{AE} = \vec 0$.
	\item We set that $\vec{OA} = \vec{OM} + \vec{MA}$. The expression $\vec{OA} + \vec{OB} + \vec{OC}$ thus
	evaluates to $3\vec{OM} + \vec{MA} + \vec{MB} + \vec{MC}$. Because $\vec{MA} + \vec{MB} + \vec{MC} = 0$,
	we see that $\vec{OM} = \frac{1}{3}(\vec{OA} + \vec{OB} + \vec{OC})$.
	\item Because $AM = \frac{2}{3} AA'$, we see that $AM + BM + CM = \frac{2}{3}(AA' + BB' + CC')$. 
	From the results of problem 7 (which we prove later) we see that $AM + BM + CM = \frac{1}{3}( 
		AB + AC + BA + BC + CA + CB
	) = 0$.
	\skipitems{1}
	\item $AB = AA' + A'B$ and $AC = AA' + A'C$, and because $A'$ is the midpoint $A'B, A'C$ cancel out. 
	Thus, $AB + AC = 2AA'$, so $AA' = \frac{1}{2}(AB + AC)$.
	\item add arrows later lmao
\end{enumerate}
\end{document}